\section{Security}\label{sec:tsps-security}

We take inspiration from the threshold signature security definitions
of~\cite{bellare2022better}, which introduce a hierarchy of increasingly
strong notions of unforgeability for fully and partially non-interactive
threshold signature schemes. A partially non-interactive scheme is one whose
$\onsign$ protocol is non-interactive, but whose message-independent
pre-processing round, analogous to $\offsign$, may be interactive; both of
our constructions from Section~\ref{sec:schemes} fall into this class.

With TS-UF-0 security, the adversary has access to a partial signature oracle
$\signo$, and the unforgeability game is defined so that the adversary cannot
win with message-signature pair $(\vec{\MOne^*}, \sigma^*)$ if the partial
signature oracle has already been queried, even once, on $\vec{\MOne^*}$.
TS-UF-1 relaxes this: the winning condition can be triggered as long as the
adversary has queried for partial signatures on $\vec{\MOne^*}$ from no more
than $t - |\Cor| - 1$ parties, where $\Cor$ is the set of corrupted parties.
When the adversary corrupts $t-1$ parties, TS-UF-0 and TS-UF-1 coincide. In
the adaptive variants, adp-TS-UF-0 and adp-TS-UF-1, the adversary additionally
has access to a corruption oracle $\sko$, alongside the partial signature
oracle $\signo$.

Let $\mathrm{TSPS} = (\setup, \keygen, \offsign, \onsign, \combine,
\verify)$ be an $(n, t)$-TSPS scheme over message space $\mathcal{M}$. The
advantage of a PPT adversary $\adv$ is defined as
%
\begin{equation*}
%
  \operatorname{Adv}_{\mathrm{TSPS}, \adv}^{\text{adp-TS-UF-1}}(\secparam) =
  \Pr\left[\mathbf{G}_{\mathrm{TS}, \adv}^{\text{adp-TS-UF-1}}(\secparam) = 1
  \right],
%
\end{equation*}
%
with game $\mathbf{G}_{\mathrm{TS}, \adv}^{\text{adp-TS-UF-1}}$ given in
Figure~\ref{fig:ts-uf-1}, and its oracles $\signo$ and $\sko$ given in
Figure~\ref{fig:oracles}. A threshold structure-preserving signature scheme
achieves adaptive TS-UF-1 security if, for every PPT adversary $\adv$,
%
\begin{equation*}
%
  \operatorname{Adv}_{\mathrm{TSPS}, \adv}^{\text{adp-TS-UF-1}}(\secparam)
  \leq \operatorname{negl}(\secparam).
%
\end{equation*}

\begin{figure}
\centering
  \fbox{\parbox{0.99\linewidth}{

      \begin{algorithmic}[1]

        \Statex \underline{$\mathbf{G}_{\mathrm{TS},\adv}^{\text
        {adp-TS-UF-1}}(\secparam):$}

        \begin{enumerate}

          \item $ \pp \from \setup (1^\secparam) $

          \item $ (n, t, \Cor) \from \adv(\pp) $

          \item $ \Hon := [1,n] \backslash \Cor $

          \item $ (\pk, \llbracket \sk \rrbracket, \llbracket \pk \rrbracket )
            \from \keygen(l, t, n) $

        \item $ (\vec{\MOne^*},\sigma^*) \from \adv^{\signo, \sko} (
            \pk, \{\sk_i\}_{i \in \Cor}, \llbracket \pk \rrbracket ) $

          \item return $ \verify (\pk, \vec{\MOne^*}, \sigma^*) \
            \wedge \ |\Cor| < t
            \wedge \ | S(\vec{\MOne^*}) | < t - | \Cor | $

      \end{enumerate}

      \end{algorithmic}

  }
}
\caption{The adp-TS-UF-1 security game for threshold structure-preserving
signature scheme TSPS, adversary $\adv$, and set $S(\vec{\MOne})$ recording
which honest signers have issued a partial signature on $\vec{\MOne}$.}
\label{fig:ts-uf-1}
\end{figure}

\begin{figure}
  \centering
  \begin{minipage}{.46\linewidth}

    \fbox{\parbox{\linewidth}{

        \begin{algorithmic}[1]

        \Statex \underline{ $ \signo (i,\vec{\MOne}): $ }

        \begin{enumerate}

          \item Assert $(\vec{\MOne} \in \mathcal{M} \wedge i \in \Hon) $

          \item $ \sigma_i \from \threshsign (\sk_i, \vec{\MOne}) $

          \item if $ \sigma_i \neq \perp $ :

            \begin{enumerate}

              \item $ S(\vec{\MOne}) \from S(\vec{\MOne}) \cup \{i\} $

            \end{enumerate}

          \item return $ \sigma_i $

        \end{enumerate}

      \end{algorithmic}
    }}

  \end{minipage}
  \hspace{.05\linewidth}
  \begin{minipage}{.46\linewidth}
    \fbox{\parbox{\linewidth}{

      \begin{algorithmic}[1]

        \Statex \underline{ $ \sko(k): $ }

        \begin{enumerate}

          \item if $ k \in \Cor $ :

            \begin{enumerate}

              \item return $ \perp $

            \end{enumerate}

          \item else : $ \Cor \leftarrow \Cor \cup \{k\} $

          \item $ \Hon \from \Hon \backslash \{k\} $

          \item return $ \sk_k $

        \end{enumerate}

      \end{algorithmic}

  }}

  \end{minipage}
  \caption{The oracles available to $\adv$ in the adp-TS-UF-1 game.}
\label{fig:oracles}
\end{figure}

\begin{theorem}\label{thm:tsps-security}

  Let $\Sigma$ be the structure-preserving signature scheme of
  Section~\ref{sec:sps-ggm} (respectively Section~\ref{sec:sps-kiltz}),
  satisfying wUF-CMA security, and let the ideal functionalities of
  Figure~\ref{fig:ideal-func} and Figure~\ref{fig:ideal-func-group} be
  realised as defined. Then the threshold structure-preserving signature
  scheme of Section~\ref{sec:tsps-abe} (respectively
  Section~\ref{sec:tsps-kiltz}) is adaptive TS-UF-1 secure.

\end{theorem}

\begin{proof}

  We argue by contrapositive. Suppose the threshold structure-preserving
  signature scheme does not achieve adaptive TS-UF-1 security. Using the
  adversary $\adv$ of the adp-TS-UF-1 game as a subroutine, we construct an
  adversary $\advTwo$ that either breaks the ideal functionalities of
  Section~\ref{sec:MPC} (case 1), or wins the wUF-CMA game of
  Figure~\ref{fig:uf-cma} against the underlying structure-preserving
  signature scheme (case 2). By design of the ideal functionalities, case 1
  does not occur for either the scheme of Section~\ref{sec:tsps-abe} or that
  of Section~\ref{sec:tsps-kiltz}, and the transformation from $\adv$ to
  $\advTwo$ that we now describe is generic to both schemes.

  $\advTwo$ forwards the setup parameters it receives to $\adv$. On receiving
  $(n, t, \Cor)$ from $\adv$, and $\pk$ from its own challenger, $\advTwo$
  samples random shares $\{\sk_i\}_{i \in [1, t-1]}$ and their corresponding
  public key shares, together with shares $\{\sk_i\}_{i \in \Cor}$ and their
  corresponding public key shares, consistent with the sharing $\llbracket
  \pk \rrbracket$ of $\pk$. It forwards $\pk$, $\{\sk_i\}_{i \in \Cor}$, and
  $\llbracket \pk \rrbracket$ to $\adv$.

  $\advTwo$ answers $\adv$'s $\signo$ queries in one of two ways. If, for the
  queried $\vec{\MOne}$, $|S(\vec{\MOne})| < t - |\Cor|$, then $\advTwo$ forms
  the requested signature share $\sigma_i$ directly from its sampled $\sk_i$;
  the information-theoretic hiding property of the linear secret sharing
  scheme means $\adv$ cannot distinguish this simulated share from a genuine
  one. Once $|S(\vec{\MOne})|$ reaches $t - |\Cor|$, $\advTwo$ instead queries
  its own signing oracle for $\vec{\MOne}$, and uses the information exchanged
  during the (possibly interactive) threshold signing protocol to construct a
  signature share $\sigma_i$ that, combined with the shares it has already
  issued, yields a signature verifying correctly under $\combine$ and
  $\verify$. In both cases, $\advTwo$ can respond to every query of $\adv$
  without ever calling its own signing oracle on a message for which it would
  later need to produce a forgery.

  Since $\advTwo$ never queries its own signing oracle on $\vec{\MOne}^*$
  itself, if $\adv$ outputs a forgery $(\vec{\MOne^*}, \sigma^*)$ satisfying
  the winning condition of the adp-TS-UF-1 game, then $(\vec{\MOne^*},
  \sigma^*)$ also satisfies the winning condition of the wUF-CMA game for
  $\advTwo$. Hence the advantage of $\adv$ in the adp-TS-UF-1 game is bounded
  by the advantage of $\advTwo$ in the wUF-CMA game of the underlying
  structure-preserving signature scheme, which is negligible by assumption.

  It remains to handle $\adv$'s queries to $\sko$. $\advTwo$ simulates these
  directly, returning the sampled $\sk_k$ for each newly corrupted party,
  until $|\Cor| = t$. Beyond this point, $\advTwo$ cannot return a further
  share $\sk_k$ consistent with $\pk$ without being able to reconstruct the
  full secret key itself, and correspondingly the adp-TS-UF-1 game's winning
  condition requires $|\Cor| < t$. For $|\Cor| < t$, every query to $\sko$ can
  be simulated as above, so the reduction goes through, and the threshold
  structure-preserving signature scheme achieves adaptive TS-UF-1 security
  with corruption bound $|\Cor| < t$.

\end{proof}
